\documentclass[12pt]{article}

\usepackage[utf8]{inputenc}
\usepackage{amsmath, amssymb, amsthm}
\usepackage{geometry}
\geometry{a4paper, margin=2.5cm}

\title{Ecuațiile Cauchy--Riemann}
\author{}
\date{}

\theoremstyle{definition}
\newtheorem{definition}{Definiție}

\theoremstyle{plain}
\newtheorem{theorem}{Teoremă}
\newtheorem{proposition}{Propoziție}

\begin{document}

\maketitle

\begin{definition}
Fie $f:\Omega \subset \mathbb{C} \to \mathbb{C}$ o funcție scrisă sub forma


\[
f(z) = u(x,y) + i\, v(x,y), \qquad z = x + i y,
\]


unde $u, v : \Omega \to \mathbb{R}$ sunt funcții reale de două variabile.
Spunem că $u$ și $v$ sunt derivabile în sens real în punctul $(x_0,y_0)$ dacă toate derivatele parțiale
$u_x, u_y, v_x, v_y$ există în acel punct.
\end{definition}

\begin{theorem}
(\textbf{Ecuațiile Cauchy--Riemann})  
Fie $f:\Omega \to \mathbb{C}$, $f(z)=u(x,y)+i\,v(x,y)$, și fie $z_0=x_0+iy_0 \in \Omega$.
Dacă $f$ este derivabilă complex în $z_0$, atunci derivatele parțiale ale lui $u$ și $v$ există în $(x_0,y_0)$ și satisfac


\[
u_x(x_0,y_0) = v_y(x_0,y_0), \qquad
u_y(x_0,y_0) = -\, v_x(x_0,y_0).
\]


Acestea se numesc \emph{ecuațiile Cauchy--Riemann}.
\end{theorem}

\begin{proposition}
(\textbf{Caracterizare punctuală})  
Dacă funcțiile $u$ și $v$ sunt derivabile în sens real într-un punct $(x_0,y_0)$ și satisfac ecuațiile Cauchy--Riemann în acel punct, atunci $f$ este derivabilă complex în $z_0 = x_0 + i y_0$.
\end{proposition}

\begin{theorem}
(\textbf{Criteriu de olomorfie})  
Fie $f:\Omega \to \mathbb{C}$, $f(z)=u(x,y)+i\,v(x,y)$, cu $u$ și $v$ funcții reale de două variabile.
Presupunem că $u$ și $v$ sunt de clasă $C^1$ pe $\Omega$ (adică au derivate parțiale continue) și că în fiecare punct $(x,y) \in \Omega$ sunt satisfăcute ecuațiile Cauchy--Riemann


\[
u_x(x,y) = v_y(x,y), \qquad
u_y(x,y) = -\, v_x(x,y).
\]


Atunci funcția $f$ este olomorfă pe $\Omega$.
\end{theorem}

\end{document}
